\chapter{Volatilité rugueuse et modèles de variance forward}\label{ch:dv:rough-volatility-and-forward-variance-models}

Prenez un long historique de volatilité réalisée quotidienne pour un indice boursier, et mesurez de combien son logarithme varie
sur un jour, deux jours, et ainsi de suite jusqu'à un mois. Tracez la moyenne des carrés de ces variations en fonction du décalage, sur des axes log-log.
Les points s'alignent sur une droite, et sa pente, divisée par deux, estime le degré de régularité de la trajectoire. Pour un modèle
de diffusion de la volatilité, tel que celui de Heston, la réponse aux petits décalages vaut un demi. En 2014, trois chercheurs
ont signalé, pour chacun des indices mesurés, une valeur voisine de 0.1~: la volatilité est beaucoup plus rugueuse que ne le supposaient les modèles utilisés par les desks.
Ce même nombre explique ce que ne peuvent pas faire les modèles du chapitre 10, à savoir
le skew à la monnaie des options sur indice, qui continue de croître à mesure que l'échéance se raccourcit. Ce chapitre part de l'objet qui rend
ces modèles praticables, la courbe de \omterm{def:dv:rough-volatility-and-forward-variance-models:fwdvar}{variance forward}. Il présente ensuite les \omterm{def:dv:rough-volatility-and-forward-variance-models:bergomi}{modèles de Bergomi} pour cette courbe, la \omterm{def:dv:rough-volatility-and-forward-variance-models:rough}{volatilité
rugueuse} et le \omterm{def:dv:rough-volatility-and-forward-variance-models:rbergomi}{modèle de Bergomi rugueux}, l'alternative dépendante du chemin, puis le test auquel doivent se soumettre tous ces modèles,
à savoir l'ajustement simultané des options sur indice et des options sur indice de volatilité.

\section{Variance forward}

Les modèles du chapitre 10 décrivent la variance instantanée $v_t$ à l'aide de quelques paramètres, et ces paramètres
déterminent la structure par terme de la volatilité à la monnaie. Les desks préfèrent prendre cette structure par terme sur le
marché et ne modéliser que sa dynamique. L'objet qui rend cela possible est la variance future
espérée, date par date.

\begin{definition}[Variance forward]\label{def:dv:rough-volatility-and-forward-variance-models:fwdvar}
La \emph{variance forward}\index{variance forward} pour la date $u$, vue à l'instant $t\le u$, est
$\xi_t(u)=\E^{\mathbb Q}_t[v_u]$, l'espérance sous la mesure de pricing de la variance instantanée en $u$. La
courbe $u\mapsto\xi_t(u)$ est la courbe de variance forward~; $\xi_0$ est celle d'aujourd'hui.
\end{definition}

La courbe se lit à partir des prix. La variance intégrée espérée jusqu'en $T$ vaut
$\E^{\mathbb Q}\bigl[\int_0^Tv_u\,du\bigr]=\int_0^T\xi_0(u)\,du$. Dans un modèle où le prix est un processus continu,
cette espérance vaut moins deux fois le prix du contrat log,
$-2\,\E^{\mathbb Q}[\ln(S_T/F_T)]$, car le lemme d'Itô donne
$\ln(S_T/F_T)=\int_0^T\sqrt{v_u}\,dW_u-\frac12\int_0^Tv_u\,du$ pour un forward sans dérive, et l'intégrale stochastique
est d'espérance nulle. Le payoff logarithmique est un strip statique d'options hors de la monnaie. Avec des taux nuls et
$k=\ln(K/F)$,
\[
\sigma_{\mathrm{VS}}^2(T)\,T=\int_0^T\xi_0(u)\,du=2\int_{-\infty}^{\infty}\mathrm{OTM}(k)\,e^{-k}\,dk,
\]
où $\mathrm{OTM}(k)$ est le prix du put en dessous du forward et du call au-dessus, par unité de forward.
Le membre de gauche est le strike équitable d'un swap de variance (le chapitre 14 construit ce produit et sa couverture), donc
$\sigma_{\mathrm{VS}}(T)$ est appelée la volatilité du swap de variance. Sa dérivée en $T$ donne la courbe~:
$\xi_0(T)=\partial_T\bigl(T\sigma_{\mathrm{VS}}^2(T)\bigr)$.

\begin{example}[La courbe de variance forward de la surface du chapitre 9]\label{ex:dv:rough-volatility-and-forward-variance-models:curve}
En intégrant le strip de chaque échéance sur la surface du chapitre 9, on obtient des volatilités de swap de variance de 16.3~\%, 18.2~\%,
20.1~\%, 22.1~\% et 23.6~\% à un mois, trois mois, six mois, un an et deux ans, contre des volatilités à la monnaie de 14.9~\%, 16.4~\%,
17.8~\%, 19.2~\% et 19.9~\%. La variance totale étant linéaire entre les piliers,
la volatilité forward $\sqrt{\xi_0}$ vaut 16.3~\% jusqu'à un mois, puis 19.1~\%, 21.8~\%, 24.0~\% et, d'un à deux ans,
24.9~\% (\cref{fig:dv:rough-volatility-and-forward-variance-models:curve}).
\end{example}

La volatilité du swap de variance se situe deux à quatre points au-dessus de la volatilité à la monnaie. Le strip pondère
chaque option par $e^{-k}$, c'est-à-dire par $1/K^2$ en strike, de sorte que les puts coûteux de bas strike
d'un smile à asymétrie négative pèsent davantage que les calls. L'écart est une mesure du skew, et il croît avec l'échéance
ici parce que les ailes s'élargissent. Une courbe de \omterm{def:dv:rough-volatility-and-forward-variance-models:fwdvar}{variance forward} doit être non négative. Un pilier où la variance totale
décroît constitue un \omterm{def:dv:implied-volatility-and-its-surface:static}{arbitrage calendaire}~: il suffit d'acheter le swap de variance de plus longue échéance et de vendre celui de plus courte. Le
\texttt{negative\_intervals} de la construction le signale avant tout calibrage du modèle.

\begin{omfigure}
\begin{tikzpicture}
\begin{axis}[width=10.5cm,height=5cm,
  xlabel={échéance (années)},ylabel={volatilité (\%)},
  tick label style={font=\small},label style={font=\small},
  xmin=0,xmax=2.05,ymin=14,ymax=26,grid=major,grid style={gray!25},
  legend columns=3,legend style={font=\footnotesize,at={(0.5,-0.3)},anchor=north,draw=none,fill=none,/tikz/every even column/.append style={column sep=8pt}},
  table/col sep=comma]
\addplot[gray,thick,mark=o,mark size=1.8pt] table[x=t,y=atm]{figdata/derivatives/12-rough-volatility-and-forward-variance-models/pillars.csv};
\addplot[omDef,very thick,mark=*,mark size=1.8pt] table[x=t,y=vs]{figdata/derivatives/12-rough-volatility-and-forward-variance-models/pillars.csv};
\addplot[omThm,very thick] table[x=t,y=fwd]{figdata/derivatives/12-rough-volatility-and-forward-variance-models/fwdvol.csv};
\legend{à la monnaie,swap de variance,{forward, $\sqrt{\xi_0(u)}$}}
\end{axis}
\end{tikzpicture}

\omcaption{La courbe de \omterm{def:dv:rough-volatility-and-forward-variance-models:fwdvar}{variance forward} de la surface du chapitre 9. La volatilité du swap de variance provient du strip d'options
de chaque échéance et se situe au-dessus de la volatilité à la monnaie parce que le strip pondère les puts sur smile.
Son forward est constant par morceaux entre les piliers. Données~: le code du chapitre.}\label{fig:dv:rough-volatility-and-forward-variance-models:curve}
\end{omfigure}

\section{Modèles de variance forward et modèle de Bergomi}

\begin{definition}[Modèle de variance forward]\label{def:dv:rough-volatility-and-forward-variance-models:fvm}
Un \emph{modèle de variance forward}\index{modèle de variance forward} spécifie la dynamique de toute la courbe
$\bigl(\xi_t(u)\bigr)_{u\ge t}$, à partir du $\xi_0$ de marché, chaque $\xi_t(u)$ étant une martingale en $t$
sous la mesure de pricing. La variance instantanée est l'extrémité courte de la courbe, $v_t=\xi_t(t)$.
\end{definition}

La condition de martingale n'est pas un choix de modélisation. $\xi_t(u)$ est une espérance conditionnelle de la même
variable aléatoire $v_u$ lorsque $t$ varie, donc elle n'a pas de dérive. Deux conséquences en découlent. Le modèle reprice chaque swap de variance
par construction, quels que soient ses paramètres, puisque la courbe initiale est une donnée. Et ses paramètres restent
libres de décrire la dynamique~: l'ampleur du mouvement de chaque point de la courbe, et la corrélation de ces mouvements
avec le sous-jacent et entre eux.

\begin{definition}[Modèle de Bergomi]\label{def:dv:rough-volatility-and-forward-variance-models:bergomi}
Le \emph{modèle de Bergomi}\index{modèle de Bergomi} est le \omterm{def:dv:rough-volatility-and-forward-variance-models:fvm}{modèle de variance forward} à dynamique log-normale et à volatilité
décroissant exponentiellement le long de la courbe. Sous sa forme à un facteur,
\[
d\xi_t(u)=\eta\,e^{-\kappa(u-t)}\,\xi_t(u)\,dW^1_t,\qquad
\xi_t(u)=\xi_0(u)\exp\Bigl(\eta e^{-\kappa(u-t)}X_t-\tfrac12\eta^2e^{-2\kappa(u-t)}\E[X_t^2]\Bigr),
\]
avec $X_t=\int_0^te^{-\kappa(t-s)}dW^1_s$ un facteur d'Ornstein--Uhlenbeck et $d\langle W^1,W^S\rangle=\rho\,dt$
pour le mouvement brownien $W^S$ du sous-jacent. La forme à deux facteurs ajoute un second facteur, à décroissance plus rapide,
et mélange les deux.
\end{definition}

Le noyau exponentiel $\eta e^{-\kappa\tau}$ est la volatilité instantanée de la \omterm{def:dv:rough-volatility-and-forward-variance-models:fwdvar}{variance forward}
$\tau=u-t$ années à l'avance. Bergomi visait à contrôler séparément trois grandeurs que Heston lie entre elles~: le
skew forward court, la \omterm{def:dv:stochastic-volatility:sv}{corrélation spot-volatilité} et la structure par terme de la \omterm{def:dv:stochastic-volatility:sv}{volatilité
de la volatilité}. Les produits dont la valeur dépend de ces grandeurs, tels que les reverse cliquets, les Napoleons et les options
sur variance réalisée, sont ceux pour lesquels il l'a conçu («~Smile dynamics II~» de Bergomi). La forme à un facteur
est markovienne en les deux variables $(S_t,X_t)$, de sorte qu'elle se simule aussi aisément que Heston. La forme à deux facteurs
s'ajuste mieux à la structure par terme observée de la \omterm{def:dv:stochastic-volatility:sv}{volatilité de la volatilité}, au prix d'un facteur supplémentaire.

Ce qu'un noyau exponentiel ne peut pas faire, c'est exploser. Quelle que soit la rapidité de $\kappa$, la volatilité d'une
\omterm{def:dv:rough-volatility-and-forward-variance-models:fwdvar}{variance forward} à un jour est au plus égale à $\eta$. Le skew de court terme du modèle est porté par la volatilité
de la \omterm{def:dv:rough-volatility-and-forward-variance-models:fwdvar}{variance forward} de court terme, il se stabilise donc comme celui de Heston (chapitre 10). Un second facteur à
décroissance rapide relève l'extrémité courte, mais la loi de puissance observée sur le marché exige un noyau infini en zéro.

\begin{omfigure}
\begin{tikzpicture}
\begin{groupplot}[group style={group size=2 by 1,horizontal sep=1.6cm},
  width=6.6cm,height=5cm,
  tick label style={font=\small},label style={font=\small},grid=major,grid style={gray!25},table/col sep=comma]
\nextgroupplot[xmode=log,ymode=log,xlabel={horizon $u-t$ (années)},ylabel={volatilité de $\xi_t(u)$},
  xmin=0.004,xmax=2,ymin=0.25,ymax=9,xtick={0.01,0.1,1},xticklabels={0.01,0.1,1},ytick={0.25,0.5,1,2,4,8},yticklabels={0.25,0.5,1,2,4,8},
  title={\small noyaux},
  legend columns=1,legend cell align=left,legend style={font=\footnotesize,at={(0.5,-0.32)},anchor=north,draw=none,fill=none}]
\addplot[omThm,very thick] table[x=tau,y=rough]{figdata/derivatives/12-rough-volatility-and-forward-variance-models/kernel.csv};
\addplot[omDef,very thick,dashed] table[x=tau,y=bergomi]{figdata/derivatives/12-rough-volatility-and-forward-variance-models/kernel.csv};
\legend{{rugueux, $\eta\sqrt{2H}\,\tau^{H-1/2}$},{Bergomi, $\eta e^{-\kappa\tau}$}}
\nextgroupplot[xlabel={jour de bourse},ylabel={volatilité instantanée (\%)},xmin=0,xmax=252,ymin=0,ymax=42,
  title={\small une année de volatilité spot},
  legend columns=2,legend style={font=\footnotesize,at={(0.5,-0.32)},anchor=north,draw=none,fill=none}]
\addplot[omThm,thin] table[x=day,y=rough]{figdata/derivatives/12-rough-volatility-and-forward-variance-models/volpaths.csv};
\addplot[omDef,very thick] table[x=day,y=bergomi]{figdata/derivatives/12-rough-volatility-and-forward-variance-models/volpaths.csv};
\legend{Bergomi rugueux,Bergomi}
\end{groupplot}
\end{tikzpicture}

\omcaption{À gauche~: la volatilité de la \omterm{def:dv:rough-volatility-and-forward-variance-models:fwdvar}{variance forward} selon l'horizon, pour le noyau rugueux ($H=0.1$, $\eta=1.9$)
et un noyau de Bergomi à un facteur calé sur lui à un mois et à un an ($\eta=2.51$, $\kappa=1.08$). Seul le noyau
rugueux explose aux horizons courts. À droite~: une année de volatilité spot quotidienne de chaque modèle, pilotée par
les mêmes accroissements browniens, avec une volatilité forward plate de 20~\%. Données~: code du
chapitre.}\label{fig:dv:rough-volatility-and-forward-variance-models:kernels}
\end{omfigure}

\section{Volatilité rugueuse}

\subsection{Les faits}

Gatheral, Jaisson et Rosenbaum ont mesuré directement la régularité de la volatilité. Ils ont pris des estimations
quotidiennes de variance réalisée pour de nombreux indices et calculé la fonction de structure
$m(q,\Delta)=\overline{|\log\sigma_{t+\Delta}-\log\sigma_t|^q}$ pour des décalages d'un à trente jours. Ils trouvent
$m(q,\Delta)\propto\Delta^{\zeta_q}$ avec $\zeta_q\approx Hq$ pour tout $q$. C'est le scaling d'un
mouvement brownien fractionnaire d'exposant de Hurst $H$ (livre~4, chapitre~17), dont les accroissements sur un
décalage $\Delta$ ont un écart type proportionnel à $\Delta^H$. Leurs estimations sont $H=0.125$ pour les contrats à terme sur le DAX, à partir de fenêtres d'une heure de données infrajournalières, et $0.142$ pour le
S\&P~500, à partir de la variance réalisée quotidienne. Parmi les 21 indices de la bibliothèque d'Oxford-Man
de variance réalisée, $\zeta_2/2$ varie de 0.075 à 0.158 (\cref{fig:dv:rough-volatility-and-forward-variance-models:rough}).
Un mouvement brownien a $H=\frac12$. Une diffusion pour la volatilité, le processus en racine carrée de Heston ou un
processus d'Ornstein--Uhlenbeck pour son logarithme, ressemble à un mouvement brownien aux décalages courts, et a donc
aussi $H=\frac12$.

\begin{definition}[Volatilité rugueuse]\label{def:dv:rough-volatility-and-forward-variance-models:rough}
La \emph{volatilité rugueuse}\index{volatilité rugueuse} est la propriété, et la classe de modèles qui la repose, selon laquelle le
logarithme de la volatilité instantanée se comporte aux courtes échelles de temps comme un mouvement brownien fractionnaire
d'exposant de Hurst $H<\frac12$, typiquement voisin de 0.1. Ses trajectoires sont plus rugueuses que celles de toute diffusion~: le
déplacement typique sur un décalage $\Delta$ varie comme $\Delta^H$ plutôt que comme $\Delta^{1/2}$.
\end{definition}

\begin{omfigure}
\begin{tikzpicture}
\begin{groupplot}[group style={group size=2 by 1,horizontal sep=1.6cm},
  width=6.6cm,height=5cm,
  tick label style={font=\small},label style={font=\small},grid=major,grid style={gray!25},table/col sep=comma]
\nextgroupplot[xmode=log,ymode=log,xlabel={décalage $\Delta$ (jours)},ylabel={$m(2,\Delta)$},
  xmin=0.9,xmax=32,xtick={1,2,5,10,20},xticklabels={1,2,5,10,20},
  title={\small fonctions de structure (simulées)},
  legend columns=1,legend cell align=left,legend style={font=\footnotesize,at={(0.5,-0.32)},anchor=north,draw=none,fill=none}]
\addplot[omThm,only marks,mark=*,mark size=1.3pt] table[x=lag,y=rough]{figdata/derivatives/12-rough-volatility-and-forward-variance-models/structure.csv};
\addplot[omDef,only marks,mark=square*,mark size=1.3pt] table[x=lag,y=ou]{figdata/derivatives/12-rough-volatility-and-forward-variance-models/structure.csv};
\addplot[omProp,only marks,mark=triangle*,mark size=1.5pt] table[x=lag,y=noisy]{figdata/derivatives/12-rough-volatility-and-forward-variance-models/structure.csv};
\legend{{rugueux, $H=0.1$~: lu 0.09},{diffusion~: lu 0.44},{diffusion bruitée~: lu 0.17}}
\nextgroupplot[xlabel={indice},ylabel={$\zeta_2/2$ publié},
  ymin=0,ymax=0.18,xmin=-0.8,xmax=20.8,ytick={0,0.05,0.1,0.15},yticklabel style={/pgf/number format/fixed,/pgf/number format/precision=2},
  xtick=data,xticklabels from table={figdata/derivatives/12-rough-volatility-and-forward-variance-models/published_h.csv}{name},
  x tick label style={rotate=90,font=\tiny,anchor=east},
  title={\small 21 indices}]
\addplot[omThm,ybar,bar width=4pt,fill=omThm!40] table[x=i,y=h]{figdata/derivatives/12-rough-volatility-and-forward-variance-models/published_h.csv};
\end{groupplot}
\end{tikzpicture}

\omcaption{À gauche~: la fonction de structure du logarithme de volatilité quotidienne simulée sur huit ans. La pente en
échelles log-log vaut $2H$. La trajectoire rugueuse a $H=0.1$~; la diffusion est un processus d'Ornstein--Uhlenbeck~; la
troisième série est la diffusion vue à travers un estimateur de variance réalisée avec une erreur de 0.08 en log. À droite~:
les estimations publiées pour 21 indices. Données~: code du chapitre~; Gatheral, Jaisson et Rosenbaum (2014),
tableau B.1.}\label{fig:dv:rough-volatility-and-forward-variance-models:rough}
\end{omfigure}

La régression est simple, et sa lecture demande du soin. La variance réalisée est elle-même une estimation. À partir de 78
rendements à cinq minutes, son erreur type relative vaut environ $\sqrt{2/78}=0.16$, de sorte que le logarithme de la
volatilité réalisée porte chaque jour une erreur voisine de 0.08. Cette erreur est indépendante d'un jour à l'autre, elle
ajoute donc une constante à $m(2,\Delta)$ à chaque décalage et aplatit la droite log-log. Sur des données simulées
(\cref{fig:dv:rough-volatility-and-forward-variance-models:rough}), une diffusion dont le $H$ vrai vaut $\frac12$, lue à travers
cette erreur, renvoie 0.17. Le moyennage sur une fenêtre joue en sens inverse~: Gatheral et ses
coauteurs ont simulé leur propre modèle rugueux avec $H=0.14$ et retrouvé de 0.16 à 0.18 à partir de
proxys de variance réalisée. La thèse de la rugosité repose donc sur plus d'une régression. L'estimation tient
d'un indice à l'autre, d'un moment $q$ à l'autre, d'un estimateur à l'autre, et elle a une conséquence dans
les prix d'options, où aucune erreur de mesure n'intervient.

\subsection{Le skew qui ne s'aplatit pas}

Définissons le skew à la monnaie $\psi(T)=\partial\sigma_{\mathrm{BS}}(k,T)/\partial k$ en $k=0$. Sur le S\&P~500
le 20 juin 2013, Gatheral, Jaisson et Rosenbaum ont ajusté $|\psi(T)|\approx A\,T^{-0.4}$ à travers les échéances.
Les modèles de volatilité stochastique diffusifs, tels que Heston, SABR et Bergomi à un facteur, donnent un skew qui tend vers une
constante lorsque $T\to0$ et décroît comme une somme d'exponentielles pour les échéances plus longues. Fukasawa a montré
qu'une volatilité pilotée par un mouvement brownien fractionnaire donne $\psi(T)\sim T^{H-1/2}$ pour les échéances courtes. Avec
$H=0.1$, l'exposant vaut $-0.4$, celui du marché. Un seul paramètre explique à la fois la rugosité de la série temporelle
et la forme du skew.

\begin{definition}[Modèle de Bergomi rugueux]\label{def:dv:rough-volatility-and-forward-variance-models:rbergomi}
Le \emph{modèle de Bergomi rugueux}\index{modèle de Bergomi rugueux} est le \omterm{def:dv:rough-volatility-and-forward-variance-models:fvm}{modèle de variance forward} à noyau en
loi de puissance. Sa variance spot vaut
\[
v_t=\xi_0(t)\exp\Bigl(\eta Y_t-\tfrac12\eta^2t^{2H}\Bigr),\qquad Y_t=\sqrt{2H}\int_0^t(t-s)^{H-1/2}\,dW^1_s,
\]
et le sous-jacent suit $dS_t/S_t=\sqrt{v_t}\,\bigl(\rho\,dW^1_t+\sqrt{1-\rho^2}\,dW^\perp_t\bigr)$. $Y$
est un mouvement brownien fractionnaire de Riemann--Liouville avec $\Var(Y_t)=t^{2H}$, donc $\E[v_t]=\xi_0(t)$~: le modèle
s'ajuste à la courbe de \omterm{def:dv:rough-volatility-and-forward-variance-models:fwdvar}{variance forward} par construction. Il possède trois paramètres, $H$, $\eta$ et $\rho$.
\end{definition}

Les \omterm{def:dv:rough-volatility-and-forward-variance-models:fwdvar}{variances forward} du modèle sont
$\xi_t(u)=\xi_0(u)\exp\bigl(\eta\sqrt{2H}\int_0^t(u-s)^{H-1/2}dW^1_s-\frac12\eta^2(u^{2H}-(u-t)^{2H})\bigr)$,
donc la volatilité de $\xi_t(u)$ vaut $\eta\sqrt{2H}\,(u-t)^{H-1/2}$. C'est le noyau en loi de puissance de la
\cref{fig:dv:rough-volatility-and-forward-variance-models:kernels}. Avec $H=0.1$ et $\eta=1.9$, il vaut 0.85 à un an,
2.30 à un mois, 4.13 à une semaine et 7.76 à un jour, soit trois fois le noyau de Bergomi apparié. La trajectoire
de la volatilité spot est en conséquence très agitée~: son logarithme varie de 0.53 par jour en moyenne, dans le panneau de droite de la
figure, contre 0.06 pour la trajectoire de Bergomi. Le prix de ce réalisme est que le modèle n'est
pas markovien. $Y_t$ dépend de tout le passé de $W^1$, avec des poids qui décroissent lentement, si bien qu'aucun
ensemble fini de variables d'état ne le résume. La méthode de Monte-Carlo (livre~4, chapitre~26) est la seule méthode générale de pricing, et une EDP
n'est pas disponible.

\begin{proposition}[Simulation exacte sur une grille]\label{prop:dv:rough-volatility-and-forward-variance-models:sim}
Sur une grille $t_1<\dots<t_n$, le vecteur $(Y_{t_1},\dots,Y_{t_n},W^1_{t_1},\dots,W^1_{t_n})$ est gaussien, avec, pour
$u\le v$,
\[
\Cov(Y_u,Y_v)=u^{2H}G(v/u),\qquad G(x)=2H\int_0^1(1-s)^{H-\frac12}(x-s)^{H-\frac12}\,ds,
\]
et, pour tous $u$ et $v$,
\[
\Cov(Y_v,W^1_u)=\frac{\sqrt{2H}}{H+\frac12}\Bigl(v^{H+\frac12}-\bigl(v-\min(u,v)\bigr)^{H+\frac12}\Bigr),
\]
si bien qu'une seule factorisation de Cholesky de cette matrice $2n\times2n$ fournit des échantillons exacts des deux processus aux dates de la grille.
\end{proposition}
\begin{proof}
Ce sont deux intégrales de Wiener contre le même $W^1$, donc le vecteur est gaussien et ses covariances sont des
intégrales de produits de noyaux. Pour $u\le v$,
$\Cov(Y_u,Y_v)=2H\int_0^u(u-s)^{H-\frac12}(v-s)^{H-\frac12}ds$~; en posant $s=us'$, on obtient $u^{2H}G(v/u)$.
Le terme croisé vaut $\sqrt{2H}\int_0^{\min(u,v)}(v-s)^{H-\frac12}ds$, qu'on intègre directement.
\end{proof}

La factorisation coûte $O(n^3)$ une fois, et chaque trajectoire $O(n^2)$. C'est peu pour les cent pas
qu'utilise le tutoriel, et coûteux pour des milliers de pas. Le schéma hybride de Bennedsen, Lunde et Pakkanen traite
le noyau exactement près de zéro, où il est singulier, et comme une fonction en escalier plus loin. La partie
lointaine est une convolution, calculée par transformée de Fourier rapide, ce qui ramène le coût à $O(n\log n)$ par
trajectoire. Les prix bénéficient ensuite d'une seconde économie, la \omterm{prop:dv:stochastic-volatility:mixing}{formule de mélange} du chapitre~10. Connaissant
la trajectoire de $W^1$, le logarithme de $S_T$ est normal, de moyenne $\rho\int\sqrt v\,dW^1-\frac12\int v\,dt$ et de
variance $(1-\rho^2)\int v\,dt$. Chaque trajectoire apporte donc un prix de Black plutôt qu'un payoff, et le bruit de
l'estimateur du skew chute fortement.

\begin{example}[La structure par terme du skew]\label{ex:dv:rough-volatility-and-forward-variance-models:skew}
Bergomi rugueux avec une volatilité forward plate de 20~\%, $H=0.1$, $\eta=1.9$ et $\rho=-0.9$ (valeurs
illustratives, non une calibration) donne un skew à la monnaie de $-1.76$ à une semaine et de $-0.32$ à un an. Une
loi de puissance ajustée d'une semaine à deux ans a un exposant de $-0.44$. Sur les mêmes échéances, la surface du
chapitre~9 donne $-1.44$ et $-0.25$, avec aussi un exposant de $-0.44$. Heston, calibré sur cette surface au chapitre~10, reste
à 0.04 de celle-ci à partir de trois mois, mais son skew à une semaine vaut $-0.70$, et son exposant sur les deux premiers mois
vaut $-0.09$~: il est presque plat (\cref{fig:dv:rough-volatility-and-forward-variance-models:skew}).
\end{example}

\begin{omfigure}
\begin{tikzpicture}
\begin{axis}[width=10.5cm,height=5.4cm,xmode=log,ymode=log,
  xlabel={échéance $T$ (années)},ylabel={skew à la monnaie $|\psi(T)|$},
  tick label style={font=\small},label style={font=\small},
  xmin=0.015,xmax=2.6,ymin=0.12,ymax=2.2,grid=major,grid style={gray!25},
  ytick={0.15,0.2,0.3,0.5,0.7,1,1.5,2},yticklabels={0.15,0.2,0.3,0.5,0.7,1,1.5,2},
  xtick={0.02,0.05,0.1,0.2,0.5,1,2},xticklabels={0.02,0.05,0.1,0.2,0.5,1,2},
  legend columns=2,legend cell align=left,legend style={font=\footnotesize,at={(0.5,-0.28)},anchor=north,draw=none,fill=none,/tikz/every even column/.append style={column sep=8pt}},
  table/col sep=comma]
\addplot[omDef,only marks,mark=*,mark size=2pt] table[x=t,y=market]{figdata/derivatives/12-rough-volatility-and-forward-variance-models/skew.csv};
\addplot[omDef,thick,dashed] table[x=t,y=fit]{figdata/derivatives/12-rough-volatility-and-forward-variance-models/skew.csv};
\addplot[omThm,very thick,mark=square*,mark size=1.6pt] table[x=t,y=rbergomi]{figdata/derivatives/12-rough-volatility-and-forward-variance-models/skew.csv};
\addplot[omProp,very thick,mark=triangle*,mark size=1.8pt] table[x=t,y=heston]{figdata/derivatives/12-rough-volatility-and-forward-variance-models/skew.csv};
\legend{surface du chapitre~9,{loi de puissance, exposant $-0.44$},{Bergomi rugueux, $H=0.1$},Heston (chapitre~10)}
\end{axis}
\end{tikzpicture}

\omcaption{Le skew à la monnaie par échéance, en axes log-log, d'une semaine à deux ans. La surface et le
Bergomi rugueux suivent des droites~; le skew de Heston se stabilise en dessous de deux mois. Données~: le tutoriel.}\label{fig:dv:rough-volatility-and-forward-variance-models:skew}
\end{omfigure}

L'ajustement du Bergomi rugueux comporte aussi une mise en garde. Son $H$ vaut 0.1, si bien que la règle des échéances courtes
$\psi\sim T^{H-1/2}$ prédit un exposant de $-0.40$. L'ajustement d'une semaine à deux ans donne $-0.44$,
parce que la règle est une limite et que deux ans ne sont pas courts. Lire $H=\frac12+\alpha$ sur un exposant ajusté
$\alpha$ est donc biaisé de quelques centièmes. L'exposant $-0.44$ de la surface se lit $H=0.06$ selon la règle, mais
environ 0.1 dès que le même ajustement sur le modèle sert d'étalon. C'est le problème du week-end de ce chapitre.

\section{Volatilité dépendante du chemin}

La rugosité est une lecture des données, parmi d'autres. Guyon et Lekeufack en ont proposé une autre en 2023. La
volatilité, selon eux, est surtout une fonction du chemin déjà parcouru par le prix.

\begin{definition}[Modèle de volatilité dépendante du chemin]\label{def:dv:rough-volatility-and-forward-variance-models:pdv}
Un \emph{modèle de volatilité dépendante du chemin}\index{modèle de volatilité dépendante du chemin} rend la volatilité
instantanée fonction des rendements passés. Dans la forme de Guyon et Lekeufack, elle s'écrit
\[
\sigma_t=\beta_0+\beta_1R_{1,t}+\beta_2\sqrt{R_{2,t}},\qquad
R_{1,t}=\int_{-\infty}^tK_1(t-s)\,\frac{dS_s}{S_s},\qquad R_{2,t}=\int_{-\infty}^tK_2(t-s)\,\Bigl(\frac{dS_s}{S_s}\Bigr)^2\!\Big/dt,
\]
une caractéristique de tendance $R_1$ (une somme pondérée des rendements passés) et une caractéristique d'activité $R_2$ (une somme pondérée des
rendements passés au carré), avec des noyaux décroissants $K_1,K_2$ dont l'intégrale vaut un et $\beta_1<0$.
\end{definition}

Guyon et Lekeufack rapportent que les rendements passés expliquent jusqu'à 90~\% de la variance de la volatilité
implicite des indices et 65~\% de celle de la volatilité réalisée. Avec des noyaux exponentiels, les caractéristiques sont markoviennes~: deux exponentielles
par noyau donnent un modèle à quatre facteurs, et Gazzani et Guyon montrent qu'il ajuste conjointement les smiles du S\&P~500 et du VIX.
Le modèle capture ce que les praticiens appellent l'effet de levier, et davantage encore~: la volatilité réagit à l'ampleur des
mouvements récents et, de façon asymétrique, à leur signe.

\begin{example}[Une grande journée]\label{ex:dv:rough-volatility-and-forward-variance-models:pdv}
Prenons $\beta_0=0.04$, $\beta_1=-0.06$, $\beta_2=0.65$ et des noyaux exponentiels de taux $\kappa_1=25$ et
$\kappa_2=15$ par an (valeurs illustratives). Sans tendance, le niveau calme est le point fixe
$\sigma^*=\beta_0/(1-\beta_2)=11.4\%$. Une journée à $-4\%$ porte la volatilité à 22.0~\%. L'écart au niveau calme est divisé par deux en dix
jours de cotation, et un mois (21 jours de cotation) plus tard la volatilité vaut encore 13.8~\%. Une journée à $+4\%$ fait d'abord baisser la volatilité, à 10.6~\%, car le terme de tendance
car le terme de tendance domine. La caractéristique d'activité décroît plus lentement, si bien que la volatilité remonte ensuite à 12.4~\%
avant de se stabiliser (\cref{fig:dv:rough-volatility-and-forward-variance-models:vix}, à gauche).
\end{example}

Les modèles rugueux et les modèles dépendants du chemin ne sont rivaux qu'en partie. Une volatilité dépendant du chemin n'a pas d'aléa
propre~: elle est pilotée par les chocs du prix, si bien que, sur un instant, la volatilité et le prix évoluent avec une corrélation $-1$,
et son noyau de tendance fait le travail que $\rho$ fait dans les autres modèles. Les différences pratiques tiennent à ce que
le modèle dépendant du chemin est markovien à peu de facteurs, et que ses paramètres viennent des séries temporelles autant
que des options.

\section{Calibration conjointe sur les options sur indice et les options sur indice de volatilité}

L'indice de volatilité (livre~1, chapitre~25) se calcule à partir d'un strip d'options sur indice à 30~jours, le même
strip qui donne la volatilité du swap de variance. Idéalisé, son carré est la \omterm{def:dv:rough-volatility-and-forward-variance-models:fwdvar}{variance forward} moyenne sur les
trente jours suivants~:
\[
\mathrm{VIX}_T^2=\frac1\Delta\int_T^{T+\Delta}\xi_T(u)\,du,\qquad\Delta=\tfrac{30}{365}.
\]
Un \omterm{def:dv:rough-volatility-and-forward-variance-models:fvm}{modèle de variance forward} price donc les contrats à terme sur le VIX et leurs options. Ce sont le test le plus direct de sa
dynamique, car ce sont des produits dont le sous-jacent est la courbe elle-même. Deux faits en découlent aussitôt. Le contrat à terme sur le VIX pour la date
$T$ vaut $\E^{\mathbb Q}[\mathrm{VIX}_T]$, alors que la volatilité forward du swap de variance vaut
$\sqrt{\E^{\mathbb Q}[\mathrm{VIX}_T^2]}$~: le contrat à terme se situe donc sous la volatilité forward, d'un terme de convexité
qui croît avec la volatilité du VIX. Par ailleurs, un modèle à noyau donné produit un seul smile de VIX, sans
plus aucun paramètre libre pour le façonner une fois que le smile de l'indice a fixé $H$, $\eta$ et $\rho$.

\begin{example}[Le VIX dans le Bergomi rugueux]\label{ex:dv:rough-volatility-and-forward-variance-models:vix}
Avec les paramètres de l'\cref{ex:dv:rough-volatility-and-forward-variance-models:skew}, le VIX à un mois a
$\E[\mathrm{VIX}^2]=0.0400$, la \omterm{def:dv:rough-volatility-and-forward-variance-models:fwdvar}{variance forward}, comme il se doit. Le contrat à terme vaut 18.74, contre une volatilité forward
du swap de variance de 20~: un écart de convexité de 1.26 point. La volatilité implicite des options sur VIX à un mois
vaut 123.1~\% pour un strike de 80~\% du contrat à terme, 123.9~\% à la monnaie et 125.7~\% à 160~\%. Le VIX du modèle
est presque lognormal, et son smile est presque plat, avec un niveau fixé par $\eta$ et $H$. Guerreiro et Guerra
font la même observation, et la mettent en regard des smiles de VIX croissants du marché.
\end{example}

\begin{omfigure}
\begin{tikzpicture}
\begin{groupplot}[group style={group size=2 by 1,horizontal sep=1.6cm},
  width=6.6cm,height=5cm,
  tick label style={font=\small},label style={font=\small},grid=major,grid style={gray!25},table/col sep=comma]
\nextgroupplot[xlabel={jours de cotation après le choc},ylabel={volatilité (\%)},xmin=-1,xmax=40,ymin=8,ymax=24,
  title={\small volatilité dépendante du chemin},
  legend columns=2,legend style={font=\footnotesize,at={(0.5,-0.32)},anchor=north,draw=none,fill=none}]
\addplot[omThm,very thick] table[x=day,y=down]{figdata/derivatives/12-rough-volatility-and-forward-variance-models/pdv.csv};
\addplot[omDef,very thick,dashed] table[x=day,y=up]{figdata/derivatives/12-rough-volatility-and-forward-variance-models/pdv.csv};
\legend{un jour à $-4\%$,un jour à $+4\%$}
\nextgroupplot[xlabel={strike du VIX (\% du contrat à terme)},ylabel={volatilité implicite (\%)},xmin=75,xmax=165,ymin=100,ymax=140,
  title={\small smile du VIX à un mois, Bergomi rugueux}]
\addplot[omThm,very thick,mark=*,mark size=1.6pt] table[x=strike,y=vol]{figdata/derivatives/12-rough-volatility-and-forward-variance-models/vix.csv};
\end{groupplot}
\end{tikzpicture}

\omcaption{À gauche~: la réponse de la volatilité dépendante du chemin de
l'\cref{ex:dv:rough-volatility-and-forward-variance-models:pdv} à une journée de forte amplitude, à la baisse ou à la hausse~; la volatilité réagit
de manière asymétrique au signe. À droite~: le smile des options sur VIX à un mois dans le Bergomi rugueux, avec les paramètres
qui donnent le skew de la \cref{fig:dv:rough-volatility-and-forward-variance-models:skew}~: presque plat. Données~: le
code du chapitre.}\label{fig:dv:rough-volatility-and-forward-variance-models:vix}
\end{omfigure}

Caler simultanément le smile du S\&P~500 et celui du VIX avec un seul modèle constitue le problème de calibration conjointe. C'est un test
sévère, car les deux marchés contraignent le même objet de deux côtés. Le smile de l'indice fixe le skew et
la volatilité de la variance à court terme, et le smile du VIX fixe directement la distribution de cette variance.
Guyon a conjecturé qu'aucun modèle à trajectoires continues ne pouvait caler les deux. Gatheral, Jusselin et Rosenbaum
ont répondu en 2020 avec le \omterm{def:dv:stochastic-volatility:heston}{modèle de Heston} rugueux quadratique, qui ajoute à la rugosité une rétroaction des rendements passés sur
la volatilité, c'est-à-dire l'ingrédient même des modèles dépendants du chemin. Guyon l'a ensuite résolu exactement par une construction non paramétrique, un problème de Schrödinger martingale qui
construit une loi conjointe calant les deux smiles, et le modèle dépendant du chemin à quatre facteurs
il les cale tous deux avec quelques paramètres. Pour un desk, la leçon est concrète. Les produits qui paient sur le VIX,
et les produits sur indice dont la valeur dépend de la volatilité de la variance, comme les cliquets et les options sur
variance réalisée, exigent un modèle testé sur les options sur VIX autant que sur les options sur indice. Un modèle calé sur le
seul smile de l'indice n'est pas testé sur la dynamique précise que ces produits rémunèrent.

\section{Tutoriel~: la volatilité rugueuse, des séries temporelles au skew}\label{tut:dv:rough-volatility-and-forward-variance-models:tut}

\begin{tutorial}
\textbf{Objectif.} Estimer la rugosité à partir d'une trajectoire de volatilité, extraire la courbe de \omterm{def:dv:rough-volatility-and-forward-variance-models:fwdvar}{variance forward} d'une surface,
simuler exactement le \omterm{def:dv:rough-volatility-and-forward-variance-models:rbergomi}{modèle de Bergomi rugueux} et caler la structure par terme du skew. \textbf{État final~:} les cinq figures du
chapitre et les résultats du problème du week-end.
\begin{tutsteps}
\item \textbf{L'estimateur de rugosité.} La fonction de structure et la pente log-log~; l'exécuter sur
\texttt{fbm(1500, 0.3, seed)} et vérifier qu'il renvoie environ 0.3~:
\omcode{code/firm/fwdvar/firm_fwdvar.py}{194}{203}{Fonction de structure et estimation de la rugosité.}\label{lst:dv:rough-volatility-and-forward-variance-models:rough}
\item \textbf{Le squelette gaussien.} La covariance conjointe du processus de Volterra et de son mouvement brownien,
à partir de la \cref{prop:dv:rough-volatility-and-forward-variance-models:sim}~:
\omcode{code/firm/fwdvar/firm_fwdvar.py}{94}{101}{Covariance de $(Y,W^1)$ sur une grille.}\label{lst:dv:rough-volatility-and-forward-variance-models:cov}
\item \textbf{Trajectoires et prix.} Simuler sur une grille avec des paires antithétiques, accumuler la somme d'Itô et la
variance intégrée, et pricer par la \omterm{prop:dv:stochastic-volatility:mixing}{formule de mélange}~:
\omcode{code/firm/fwdvar/firm_fwdvar.py}{114}{143}{Bergomi rugueux par simulation exacte et formule de mélange.}\label{lst:dv:rough-volatility-and-forward-variance-models:rbergomi}
\item \textbf{Exécuter} \texttt{dv\_roughvol.roughness\_table()}, \texttt{variance\_curve()}, \texttt{skew\_term()},
\texttt{vix\_smile()} et \texttt{fig\_roughvol.py}.
\end{tutsteps}
\textbf{Pour aller plus loin.} Remplacer les 100 pas fixes par 25 et 400 et observer le skew à une semaine~; poser $H=0.3$ et
recaler l'exposant~; ajouter un bruit de 0.04 au lieu de 0.08 au logarithme de la volatilité de la diffusion et réestimer $H$.
\end{tutorial}

\section{Construction~: variance forward et Bergomi rugueux}\label{bld:dv:rough-volatility-and-forward-variance-models:fwdvar}

\begin{build}
\bfield{Rôle} La courbe de \omterm{def:dv:rough-volatility-and-forward-variance-models:fwdvar}{variance forward} de la firme miniature, que lit chaque produit de volatilité de la partie~III
(swaps de variance, contrats à terme sur le VIX, cliquets), et un moteur de \omterm{def:dv:rough-volatility-and-forward-variance-models:rough}{volatilité rugueuse} pour les options sur indice à court terme.
\bfield{Interface} \texttt{log\_strip\_variance(vol\_of\_k, t)}~; \texttt{ForwardVarianceCurve}\allowbreak\texttt{.from\_vs\_vols(times, vols)} avec \texttt{xi(u)}, \texttt{total\_variance(t)}, \texttt{vs\_vol(t)}, \texttt{forward\_vol(t1, t2)},
\texttt{negative\_intervals()}~; \texttt{joint\_cov(times, H)}~; \texttt{rbergomi\_integrals(xi0, H, eta, t,
steps, n\_paths, seed)}~; \texttt{mixing\_calls}~; \texttt{rbergomi\_smile}~; \texttt{rbergomi\_vix}~; \texttt{atm\_skew}~;
\texttt{power\_law\_fit}~; \texttt{fbm}~; \texttt{structure\_function}~; \texttt{roughness}.
\bfield{Règles} La variance totale est non décroissante, sinon la courbe est rejetée avant que tout modèle ne la voie~;
$\E[v_t]=\xi_0(t)$ exactement dans le simulateur (le compensateur de la variance utilise le même $t^{2H}$ que la
covariance)~; les prix simulés sont vérifiés par test de martingale (le forward est recalculé à l'intérieur de son erreur type)~;
les graines sont explicites.
\bfield{Tests d'acceptation} \texttt{code/firm/fwdvar/tests/}~: le strip d'un smile plat restitue sa variance~; les morceaux de la courbe
s'additionnent~; $G(1)=1$ et la covariance se réduit à celle du mouvement brownien en $H=\frac12$~; avec $\eta=0$ et
$\rho=0$, le smile est plat à $\sqrt{\xi_0}$~; une corrélation négative donne un skew négatif~;
$\E[\mathrm{VIX}^2]=\xi_0$~; l'estimateur retrouve $H$ à partir de trajectoires exactes de mouvement brownien fractionnaire.
\bfield{Pour aller plus loin} Le schéma hybride avec une transformée de Fourier rapide, pour des grilles de milliers de pas~; un
\omterm{def:dv:rough-volatility-and-forward-variance-models:bergomi}{modèle de Bergomi} à deux facteurs avec la même interface~; le modèle dépendant du chemin à quatre facteurs et son VIX.
\end{build}

\begin{omsources}
\item J.~Gatheral, T.~Jaisson et M.~Rosenbaum, «~Volatility is rough~», \emph{Quantitative Finance} 18(6) (2018)
933--949; arXiv 1410.3394.
\item C.~Bayer, P.~Friz et J.~Gatheral, «~Pricing under rough volatility~», \emph{Quantitative Finance} 16(6)
(2016) 887--904.
\item M.~Bennedsen, A.~Lunde et M.~S.~Pakkanen, «~Hybrid scheme for Brownian semistationary processes~»,
\emph{Finance and Stochastics} 21(4) (2017) 931--965.
\item L.~Bergomi, «~Smile dynamics II~», SSRN 1493302.
\item J.~Guyon et J.~Lekeufack, «~Volatility is (mostly) path-dependent~», \emph{Quantitative Finance} 23(9)
(2023) 1221--1258.
\item G.~Gazzani et J.~Guyon, «~Pricing and calibration in the 4-factor path-dependent volatility model~»,
arXiv 2406.02319 (2024).
\item J.~Gatheral, P.~Jusselin et M.~Rosenbaum, «~The quadratic rough Heston model and the joint S\&P~500/VIX
smile calibration problem~», arXiv 2001.01789 (2020).
\item J.~Guyon, «~Dispersion-constrained martingale Schrödinger problems and the exact joint S\&P~500/VIX smile
calibration puzzle~», \emph{Finance and Stochastics} (2023).
\end{omsources}

\section{Exercices}

\begin{exercise}[$\star$]\label{exo:dv:rough-volatility-and-forward-variance-models:1}
Les volatilités de swap de variance à un an et à deux ans valent 20~\% et 22~\%. Calculer la volatilité forward entre
un an et deux ans.
\end{exercise}

\begin{exercise}[$\star$]\label{exo:dv:rough-volatility-and-forward-variance-models:2}
Pourquoi la volatilité du swap de variance est-elle supérieure à la volatilité à la monnaie lorsque le skew est orienté vers la baisse~?
\end{exercise}

\begin{exercise}[$\star$]\label{exo:dv:rough-volatility-and-forward-variance-models:3}
Doubler le retard multiplie $m(2,\Delta)$ par 1.15. Estimer $H$.
\end{exercise}

\begin{exercise}[$\star\star$]\label{exo:dv:rough-volatility-and-forward-variance-models:4}
Montrer que, dans le \omterm{def:dv:rough-volatility-and-forward-variance-models:rbergomi}{modèle de Bergomi rugueux}, $\E[v_t]=\xi_0(t)$ et que la volatilité de $\xi_t(u)$ vaut
$\eta\sqrt{2H}\,(u-t)^{H-1/2}$.
\end{exercise}

\begin{exercise}[$\star\star$]\label{exo:dv:rough-volatility-and-forward-variance-models:5}
Dans l'\cref{ex:dv:rough-volatility-and-forward-variance-models:vix}, le contrat à terme sur le VIX vaut 18.74 alors que la volatilité du
swap de \omterm{def:dv:rough-volatility-and-forward-variance-models:fwdvar}{variance forward} vaut 20. Expliquer l'écart et dire ce qu'une valeur plus grande de $\eta$ ferait à cet écart.
\end{exercise}

\begin{exercise}[$\star\star$]\label{exo:dv:rough-volatility-and-forward-variance-models:6}
Dans l'\cref{ex:dv:rough-volatility-and-forward-variance-models:pdv}, dériver le niveau calme $\sigma^*$ et expliquer pourquoi
la volatilité baisse d'abord après une journée à $+4\%$, puis remonte au-dessus de $\sigma^*$.
\end{exercise}

\begin{exercise}[$\star\star\star$]\label{exo:dv:rough-volatility-and-forward-variance-models:7}
\emph{Programmation.} Estimer $H$ sur la trajectoire de diffusion du chapitre (log-volatilité de type Ornstein--Uhlenbeck) sans bruit et
avec un bruit de variance réalisée de 0.08 en log, sur cinq graines. Quelle plage obtient-on~?
\end{exercise}

\begin{exercise}[$\star\star\star$]\label{exo:dv:rough-volatility-and-forward-variance-models:8}
\emph{Trouver l'erreur.} «~La fonction de structure de notre volatilité réalisée quotidienne donne $H=0.17$, très inférieur à un demi,
donc le modèle de volatilité diffusive du desk est réfuté.~»
\end{exercise}

\section{Problème~: Le skew qui refuse de s'aplatir}

\begin{problem}[{Problème du week-end --- la loi de puissance du skew à la monnaie}]\label{pb:dv:rough-volatility-and-forward-variance-models:1}
Un desk d'indices price des options à courte échéance avec le \omterm{def:dv:stochastic-volatility:heston}{modèle de Heston} du chapitre 10 et veut savoir si un modèle
rugueux le servirait mieux. Il dispose de la surface du chapitre 9, de la calibration de Heston et du
simulateur de Bergomi rugueux du chapitre.

\medskip\noindent\textbf{Partie I --- La courbe.}
\begin{enumerate}
\item Donner les volatilités du swap de variance à un an et à la monnaie de la surface.
\item Donner la volatilité forward entre six mois et un an.
\item Si la volatilité du swap de variance à un an augmente d'un point et que celle à six mois ne bouge pas, que
devient cette volatilité forward~?
\item Donner le strike en volatilité d'un swap de variance qui débute dans un an et se termine dans deux ans.
\item La courbe est-elle exempte d'\omterm{def:dv:implied-volatility-and-its-surface:static}{arbitrage calendaire}~?
\end{enumerate}

\medskip\noindent\textbf{Partie II --- Noyaux.}
\begin{enumerate}[resume]
\item Donner la volatilité de la \omterm{def:dv:rough-volatility-and-forward-variance-models:fwdvar}{variance forward} du noyau rugueux à une semaine, à un mois et à un an ($H=0.1$,
$\eta=1.9$).
\item Donner le noyau de Bergomi à un facteur qui le reproduit à un mois et à un an.
\item De quel facteur le noyau rugueux dépasse-t-il celui-ci un jour à l'avance~?
\item Pourquoi aucun choix de $\kappa$ ne permet-il au noyau exponentiel de reproduire le noyau rugueux à toutes les échéances~?
\item Que ferait un second facteur de Bergomi, rapide~?
\end{enumerate}

\medskip\noindent\textbf{Partie III --- Skews.}
\begin{enumerate}[resume]
\item Donner le skew à la monnaie de la surface à une semaine et à un an.
\item Donner celui de Heston aux mêmes échéances.
\item Donner les exposants de loi de puissance ajustés d'une semaine à deux ans pour la surface et pour le Bergomi rugueux.
\item Donner l'exposant de Heston sur les deux premiers mois.
\item Confronter l'exposant de la surface au rapport de ses skews à une semaine et à un an.
\end{enumerate}

\medskip\noindent\textbf{Partie IV --- Jugement.}
\begin{enumerate}[resume]
\item Quelle valeur de $H$ la règle des courtes échéances $\psi\sim T^{H-1/2}$ lit-elle dans l'exposant de la surface~?
\item Le Bergomi rugueux avec $H=0.1$ donne le même exposant ajusté. Qu'est-ce que cela dit de cette lecture~?
\item Comparer avec le $H$ estimé à partir de séries temporelles d'indices.
\item Énoncer le \emph{résultat nommé}~: l'exposant de la loi de puissance de la structure par terme du skew à la monnaie ajusté sur la
surface du chapitre, et l'exposant de Hurst qui en découle.
\item Faut-il transférer la tarification à courte échéance du desk vers le \omterm{def:dv:rough-volatility-and-forward-variance-models:rbergomi}{modèle de Bergomi rugueux}~? Nommer un coût.
\end{enumerate}
\end{problem}

\section{Questions d'entretien}

\begin{interviewq}[$\star$ \iqroles{trader, chercheur}]\label{iq:dv:rough-volatility-and-forward-variance-models:1}
Qu'est-ce que la \omterm{def:dv:rough-volatility-and-forward-variance-models:fwdvar}{variance forward}, et comment la lisez-vous sur le marché~?
\end{interviewq}

\begin{interviewq}[$\star$ \iqroles{chercheur}]\label{iq:dv:rough-volatility-and-forward-variance-models:2}
Que signifie «~la volatilité est rugueuse~», et quelles en sont les preuves~?
\end{interviewq}

\begin{interviewq}[$\star\star$ \iqroles{chercheur, trader}]\label{iq:dv:rough-volatility-and-forward-variance-models:3}
Pourquoi le \omterm{def:dv:stochastic-volatility:heston}{modèle de Heston} ne parvient-il pas à ajuster le skew à courte échéance, et comment la \omterm{def:dv:rough-volatility-and-forward-variance-models:rough}{volatilité rugueuse} y remédie-t-elle~?
\end{interviewq}

\begin{interviewq}[$\star\star$ \iqroles{développeur}]\label{iq:dv:rough-volatility-and-forward-variance-models:4}
Comment simuleriez-vous le \omterm{def:dv:rough-volatility-and-forward-variance-models:rbergomi}{modèle de Bergomi rugueux}, et quel en est le coût~?
\end{interviewq}

\begin{interviewq}[$\star\star$ \iqroles{chercheur, risque}]\label{iq:dv:rough-volatility-and-forward-variance-models:5}
Qu'est-ce que le problème de calibration conjointe S\&P~500/VIX, et pourquoi un desk s'en soucie-t-il~?
\end{interviewq}

\begin{interviewq}[$\star\star\star$ \iqroles{trader, risque}]\label{iq:dv:rough-volatility-and-forward-variance-models:6}
Deux modèles ajustent aussi bien la surface indicielle d'aujourd'hui~: le \omterm{def:dv:rough-volatility-and-forward-variance-models:bergomi}{modèle de Bergomi} à un facteur et le \omterm{def:dv:rough-volatility-and-forward-variance-models:rbergomi}{modèle de Bergomi rugueux}. Lesquels de vos produits
ces deux modèles pricent-ils différemment, et comment trancheriez-vous~?
\end{interviewq}
